% 量子力学1 量子力学1-10 問題3（037-real。当時は 量子力学1-5 問題3） 解いた日 2026-10-30
%% kotae
\toi{1}{◎}{}
\[ I_{1}=\int\psi^{*}(\hat{p}\psi)\dd x \]
\begin{align*}
\hat{p}\psi&=-\ii\hbar\frac{\dd}{\dd x}\Bigl((\pi b^{2})^{-1/4}\ee^{-x^{2}/(2b^{2})}\Bigr)\\
&=-\ii\hbar(\pi b^{2})^{-1/4}\cdot\Bigl(-\frac{x}{b^{2}}\Bigr)\ee^{-x^{2}/(2b^{2})}\\
&=\ii\hbar(\pi b^{2})^{-1/4}\cdot\frac{x}{b^{2}}\ee^{-x^{2}/(2b^{2})}
\end{align*}
\begin{align*}
\psi^{*}\hat{p}\psi&=(\pi b^{2})^{-1/4}\ee^{-x^{2}/(2b^{2})}\\
&\qquad\cdot\ii\hbar(\pi b^{2})^{-1/4}\frac{x}{b^{2}}\ee^{-x^{2}/(2b^{2})}\\
&=\ii\hbar(\pi b^{2})^{-1/2}\frac{x}{b^{2}}\ee^{-x^{2}/b^{2}}
\end{align*}
\begin{align*}
I_{1}&=-\frac{1}{2}\ii\hbar(\pi b^{2})^{-1/2}\int\frac{-2x}{b^{2}}\ee^{-x^{2}/b^{2}}\dd x\\
&=\kakimachigai{\frac{1}{2}}{-\frac{1}{2}}\ii\hbar(\pi b^{2})^{-1/2}\Bigl[\ee^{-x^{2}/b^{2}}\Bigr]_{-\infty}^{\infty}\\
&=0
\end{align*}
\begin{align*}
I_{2}&=\int\psi^{*}x\Bigl(-\ii\hbar\frac{\dd}{\dd x}(\psi)\Bigr)\dd x\\
&=\int\ii\hbar(\pi b^{2})^{-1/2}\frac{x^{2}}{b^{2}}\ee^{-x^{2}/b^{2}}\dd x\\
&=\ii\hbar(\pi b^{2})^{-1/2}\cdot\Bigl(-\frac{1}{2}\Bigr)\int x\cdot\bigl(\ee^{-x^{2}/b^{2}}\bigr)'\dd x\\
&=-\frac{\ii\hbar}{2}(\pi b^{2})^{-1/2}\biggl(\Bigl[x\ee^{-x^{2}/b^{2}}\Bigr]_{-\infty}^{\infty}\\
&\qquad-\int\ee^{-x^{2}/b^{2}}\dd x\biggr)\\
&=\frac{\ii\hbar}{2}\frac{1}{\sqrt{\pi b^{2}}}\cdot\sqrt{\pi b^{2}}\\
&=\frac{\ii\hbar}{2}
\end{align*}
\[ I_{3}=\int\psi^{*}\Bigl(-\ii\hbar\frac{\dd}{\dd x}(x\psi)\Bigr)\dd x \]
\begin{align*}
\frac{\dd}{\dd x}(x\psi)&=\psi+\kakimachigai{x\psi^{*}}{x\psi'}\\
&=(\pi b^{2})^{-1/4}\ee^{-x^{2}/(2b^{2})}\\
&\qquad-\frac{x^{2}}{b^{2}}(\pi b^{2})^{-1/4}\ee^{-x^{2}/(2b^{2})}\\
&=(\pi b^{2})^{-1/4}\Bigl(1-\frac{x^{2}}{b^{2}}\Bigr)\ee^{-x^{2}/(2b^{2})}
\end{align*}
\begin{align*}
&\psi^{*}\Bigl(-\ii\hbar\frac{\dd}{\dd x}(x\psi)\Bigr)\\
&=-\ii\hbar(\pi b^{2})^{-1/2}\Bigl(1-\frac{x^{2}}{b^{2}}\Bigr)\ee^{-x^{2}/b^{2}}
\end{align*}
\begin{align*}
I_{3}&=\ii\hbar(\pi b^{2})^{-1/2}\\
&\qquad\cdot\int\Bigl(\frac{x^{2}}{b^{2}}\ee^{-x^{2}/b^{2}}-\ee^{-x^{2}/b^{2}}\Bigr)\dd x
\end{align*}
\[ \int\frac{x^{2}}{b^{2}}\ee^{-x^{2}/b^{2}}\dd x=\frac{1}{2}\sqrt{\pi b^{2}} \]
\begin{align*}
I_{3}&=\ii\hbar(\pi b^{2})^{-1/2}\Bigl(\frac{1}{2}\sqrt{\pi b^{2}}-\sqrt{\pi b^{2}}\Bigr)\\
&=-\frac{\ii\hbar}{2}
\end{align*}
\[ I_{4}=\frac{I_{2}+I_{3}}{2}=0 \]
\toi{2}{◎}{}
\[ (\hat{x}\hat{p})^{\dagger}=\hat{p}^{\dagger}\hat{x}^{\dagger}=\hat{p}\hat{x}\neq\hat{x}\hat{p} \]
エルミートでない。
\begin{align*}
\Bigl(\frac{\hat{x}\hat{p}+\hat{p}\hat{x}}{2}\Bigr)^{\dagger}
&=\frac{\hat{p}^{\dagger}\hat{x}^{\dagger}+\hat{x}^{\dagger}\hat{p}^{\dagger}}{2}\\
&=\frac{\hat{x}\hat{p}+\hat{p}\hat{x}}{2}
\end{align*}
よりエルミート。

$\hat{x}$，$\hat{p}$ はそれぞれエルミートだが，$(\hat{x}\hat{p})$ はエルミートでないため矛盾しない。
\toi{3}{◎}{}
\begin{align*}
I_{2}-I_{3}&=\int\bigl(\psi^{*}\hat{x}\hat{p}\psi\kakimachigai{+}{-}\psi^{*}\hat{p}\hat{x}\psi\bigr)\dd x\\
&=\int\psi^{*}\comm{\hat{x}}{\hat{p}}\psi\dd x\\
&=\ii\hbar\int\abs{\psi}^{2}\dd x\\
&=\ii\hbar
\end{align*}
\[ I_{2}-I_{3}=\frac{\ii\hbar}{2}-\Bigl(-\frac{\ii\hbar}{2}\Bigr)=\ii\hbar \]
%% kizuki
